\documentclass[25pt,a0paper,portrait]{tikzposter}

\usepackage{fontspec}
\usepackage{polyglossia}

\setdefaultlanguage{russian}
\setotherlanguage{english}

\setmainfont{PT Serif}
\setsansfont{PT Sans}
\setmonofont{PT Mono}

\newfontfamily\cyrillicfont{PT Serif}[Script=Cyrillic]
\newfontfamily\cyrillicfontsf{PT Sans}[Script=Cyrillic]
\newfontfamily\cyrillicfonttt{PT Mono}[Script=Cyrillic]

\usepackage{graphicx}
\usepackage{booktabs}

\usepackage[
    backend=biber,
    style=numeric
]{biblatex}

\addbibresource{references.bib}

\title{Пример научного постера}
\author{Елизавета Гайдамака}
\institute{Российский университет дружбы народов}

\usetheme{Basic}

\begin{document}

\maketitle

\begin{columns}

% Левая колонка
\column{0.5}

\block{Введение}{
    В данной работе рассматривается пример оформления
    научного постера средствами LaTeX.

    LaTeX позволяет создавать воспроизводимые документы
    и автоматически оформлять библиографические ссылки
    \cite{lamport1994}.
}

\block{Цель работы}{
    Цель работы — изучить возможности пакета
    \texttt{tikzposter} и создать научный постер.
}

\block{Изображение}{
    \begin{center}
        \includegraphics[width=0.8\linewidth]{image/1.png}

        Рис. 1. Пример изображения
    \end{center}
}

% Правая колонка
\column{0.5}

\block{Результаты}{

    Основные результаты представлены в таблице.

    \begin{center}
    \begin{tabular}{lcc}
        \toprule
        Метод & Время, с & Точность, \% \\
        \midrule
        Метод A & 12.4 & 91 \\
        Метод B & 10.8 & 94 \\
        Метод C & 9.7  & 96 \\
        \bottomrule
    \end{tabular}
    \end{center}
}

\block{Обсуждение}{
    Полученные результаты показывают, что использование
    LaTeX удобно для подготовки научных материалов.

    Пакет \texttt{tikzposter} позволяет использовать
    колонки, блоки, изображения и таблицы.
}

\block{Выводы}{
    В ходе работы был создан научный постер,
    содержащий основные элементы научного представления
    результатов.
}

\block{Библиография}{
    \printbibliography[heading=none]
}

\end{columns}

\end{document}